% fig16_data_flywheel.tex — 图16 大模型技术栈运行闭环（数据飞轮）
% 《深度学习与大语言模型：前沿技术全景报告》图示库
% 由 dl_llm_report.tex 抽取，经 % fig16_data_flywheel.tex — 图16 大模型技术栈运行闭环（数据飞轮）
% 《深度学习与大语言模型：前沿技术全景报告》图示库
% 由 dl_llm_report.tex 抽取，经 % fig16_data_flywheel.tex — 图16 大模型技术栈运行闭环（数据飞轮）
% 《深度学习与大语言模型：前沿技术全景报告》图示库
% 由 dl_llm_report.tex 抽取，经 % fig16_data_flywheel.tex — 图16 大模型技术栈运行闭环（数据飞轮）
% 《深度学习与大语言模型：前沿技术全景报告》图示库
% 由 dl_llm_report.tex 抽取，经 \input{figs/fig16_data_flywheel} 引用（caption/label 在主文件）
% 注意：本文件为片段，依赖主文件导言区的 tikz/pgfplots 宏包、颜色定义与 tikzset 样式，不可单独编译。

\begin{tikzpicture}
\node[blk, minimum width=31mm] (app) at (0,2.4) {应用 / 用户};
\node[grn, minimum width=31mm] (data) at (5.4,2.4) {数据工程\\ 交互 · 合成 · 标注};
\node[blk, minimum width=31mm] (train) at (10.8,2.4) {训练\\ 预训练 / SFT / RL};
\node[orn, minimum width=31mm] (deploy) at (10.8,-0.6) {部署\\ vLLM · 端侧 · Agent};
\node[gry, minimum width=31mm] (ev) at (0,-0.6) {评测与安全\\ 基准 · 红队 · 备案};
\draw[brr] (app) -- node[lab, above=1mm] {使用} (data);
\draw[brr] (data) -- node[lab, above=1mm] {语料 / 偏好} (train);
\draw[arr] (train) -- node[lab, right=1mm] {权重} (deploy);
\draw[arr] (deploy) -- node[lab, above=1mm] {线上指标} (ev);
\draw[arr] (ev) -- node[lab, left=1mm] {失效模式回炉} (app);
\node[lab] at (5.4,0.9) {—— 数据飞轮：每一轮迭代为下一轮积累数据与评估资产 ——};
\end{tikzpicture}
 引用（caption/label 在主文件）
% 注意：本文件为片段，依赖主文件导言区的 tikz/pgfplots 宏包、颜色定义与 tikzset 样式，不可单独编译。

\begin{tikzpicture}
\node[blk, minimum width=31mm] (app) at (0,2.4) {应用 / 用户};
\node[grn, minimum width=31mm] (data) at (5.4,2.4) {数据工程\\ 交互 · 合成 · 标注};
\node[blk, minimum width=31mm] (train) at (10.8,2.4) {训练\\ 预训练 / SFT / RL};
\node[orn, minimum width=31mm] (deploy) at (10.8,-0.6) {部署\\ vLLM · 端侧 · Agent};
\node[gry, minimum width=31mm] (ev) at (0,-0.6) {评测与安全\\ 基准 · 红队 · 备案};
\draw[brr] (app) -- node[lab, above=1mm] {使用} (data);
\draw[brr] (data) -- node[lab, above=1mm] {语料 / 偏好} (train);
\draw[arr] (train) -- node[lab, right=1mm] {权重} (deploy);
\draw[arr] (deploy) -- node[lab, above=1mm] {线上指标} (ev);
\draw[arr] (ev) -- node[lab, left=1mm] {失效模式回炉} (app);
\node[lab] at (5.4,0.9) {—— 数据飞轮：每一轮迭代为下一轮积累数据与评估资产 ——};
\end{tikzpicture}
 引用（caption/label 在主文件）
% 注意：本文件为片段，依赖主文件导言区的 tikz/pgfplots 宏包、颜色定义与 tikzset 样式，不可单独编译。

\begin{tikzpicture}
\node[blk, minimum width=31mm] (app) at (0,2.4) {应用 / 用户};
\node[grn, minimum width=31mm] (data) at (5.4,2.4) {数据工程\\ 交互 · 合成 · 标注};
\node[blk, minimum width=31mm] (train) at (10.8,2.4) {训练\\ 预训练 / SFT / RL};
\node[orn, minimum width=31mm] (deploy) at (10.8,-0.6) {部署\\ vLLM · 端侧 · Agent};
\node[gry, minimum width=31mm] (ev) at (0,-0.6) {评测与安全\\ 基准 · 红队 · 备案};
\draw[brr] (app) -- node[lab, above=1mm] {使用} (data);
\draw[brr] (data) -- node[lab, above=1mm] {语料 / 偏好} (train);
\draw[arr] (train) -- node[lab, right=1mm] {权重} (deploy);
\draw[arr] (deploy) -- node[lab, above=1mm] {线上指标} (ev);
\draw[arr] (ev) -- node[lab, left=1mm] {失效模式回炉} (app);
\node[lab] at (5.4,0.9) {—— 数据飞轮：每一轮迭代为下一轮积累数据与评估资产 ——};
\end{tikzpicture}
 引用（caption/label 在主文件）
% 注意：本文件为片段，依赖主文件导言区的 tikz/pgfplots 宏包、颜色定义与 tikzset 样式，不可单独编译。

\begin{tikzpicture}
\node[blk, minimum width=31mm] (app) at (0,2.4) {应用 / 用户};
\node[grn, minimum width=31mm] (data) at (5.4,2.4) {数据工程\\ 交互 · 合成 · 标注};
\node[blk, minimum width=31mm] (train) at (10.8,2.4) {训练\\ 预训练 / SFT / RL};
\node[orn, minimum width=31mm] (deploy) at (10.8,-0.6) {部署\\ vLLM · 端侧 · Agent};
\node[gry, minimum width=31mm] (ev) at (0,-0.6) {评测与安全\\ 基准 · 红队 · 备案};
\draw[brr] (app) -- node[lab, above=1mm] {使用} (data);
\draw[brr] (data) -- node[lab, above=1mm] {语料 / 偏好} (train);
\draw[arr] (train) -- node[lab, right=1mm] {权重} (deploy);
\draw[arr] (deploy) -- node[lab, above=1mm] {线上指标} (ev);
\draw[arr] (ev) -- node[lab, left=1mm] {失效模式回炉} (app);
\node[lab] at (5.4,0.9) {—— 数据飞轮：每一轮迭代为下一轮积累数据与评估资产 ——};
\end{tikzpicture}
